% ══════════════════════════════════════════════════════════════════════════════
% Lecture 14: Strongly Connected Components and Flow Networks
% ══════════════════════════════════════════════════════════════════════════════

% ══════════════════════════════════════════════════════════════════════════════
\section*{Strongly Connected Components (SCC)}
\addcontentsline{toc}{subsection}{Strongly Connected Components}
% ══════════════════════════════════════════════════════════════════════════════

% ─────────────────────────────────────────────────────────────────────────────
\subsection{Definition}

\begin{definition}{Strongly Connected Component}{scc}
A \textbf{strongly connected component (SCC)} of a directed graph
$G = (V, E)$ is a \emph{maximal} set of vertices $C \subseteq V$
such that for every $u, v \in C$, there is a path from $u$ to $v$
\emph{and} a path from $v$ to $u$ (denoted $u \leadsto v$ and $v \leadsto u$).
\end{definition}

\textbf{Example:} on a graph with 10 vertices $\{a, b, c, d, e, f, g, h, i, j\}$,
the SCCs may group as $\{a, b, e, f\}$, $\{c, d, h\}$, $\{g\}$, $\{i, j\}$
depending on the edges present. A subset that is not maximal, or in which
two vertices are not mutually reachable, is not an SCC.

% ─────────────────────────────────────────────────────────────────────────────
\subsection{Component Graph}

\begin{definition}{Component Graph $G^{SCC}$}{gscc}
For a directed graph $G = (V, E)$, the component graph
$G^{SCC} = (V^{SCC}, E^{SCC})$ is defined by:
\begin{itemize}[noitemsep]
  \item $V^{SCC}$ contains one vertex for each SCC of $G$.
  \item $E^{SCC}$ contains an edge between two SCCs if there exists an edge
        between the corresponding SCCs in $G$.
\end{itemize}
\end{definition}

\begin{theorem}{$G^{SCC}$ is a DAG}{sccdag}
The component graph $G^{SCC}$ is always a directed acyclic graph.
\end{theorem}

This follows from the fact that if a cycle existed in $G^{SCC}$, all involved SCCs
would form a single SCC, which contradicts their definition.

% ─────────────────────────────────────────────────────────────────────────────
\subsection{Kosaraju's Algorithm}

\begin{algorithm}[H]
\caption{SCC($G$)}
\begin{algorithmic}[1]
\State Call \textsc{DFS}($G$) to compute finishing times $u.f$ for all $u$
\State Compute the transpose graph $G^T$   \AlgComment{reverse all edges}
\State Call \textsc{DFS}($G^T$) processing vertices in
       \textbf{decreasing} order of $u.f$ calculated in step 1
\State \Return each tree of the DFS forest from step 3 as a distinct SCC
\end{algorithmic}
\end{algorithm}

\textbf{Transpose Graph $G^T$:}
$G^T = (V, E^T)$ where $E^T = \{(u,v) : (v,u) \in E\}$ — all edges are
reversed. $G^T$ is calculated in $\Theta(V + E)$ using adjacency lists, and
$G$ and $G^T$ possess exactly the same SCCs.

\begin{complexity}{Kosaraju's SCC}{scccplx}
Each step executes in $\Theta(V + E)$, so the
\textbf{total complexity is $\Theta(V + E)$}.
\end{complexity}

\subsubsection{Intuition of Correctness}

The first DFS orders the SCCs in topological order in $G^{SCC}$ (which is a DAG).
The second DFS on $G^T$ starts from the SCC with the largest $u.f$, and in $G^T$ the
edges between SCCs are reversed: one cannot therefore "escape" to another
SCC. Each tree of the second DFS is exactly one SCC.

\subsubsection{Example}

On a graph with 10 vertices, the first DFS (on $G$) produces the
following timestamps (reading the slides: $d$: 1/12, $i$: 10/11, $h$: 2/5, $c$: 3/4,
$j$: 7/8, $e$: 6/9, $a$: 14/19, $b$: 15/16, $f$: 17/18, $g$: 13/20).

The second DFS on $G^T$ starts from $g$ ($f=20$), then $a$ ($f=19$), etc., and
each tree formed corresponds to an SCC:

\begin{center}
\begin{tikzpicture}[font=\small, >=Stealth,
    scc/.style={draw, rounded corners, fill=blue!10,
                minimum width=1.4cm, minimum height=0.6cm, inner sep=3pt}]
  \node[scc, fill=green!15]  (s1) at (0, 0)   {$\{a,b,f,g\}$};
  \node[scc, fill=orange!15] (s2) at (3, 0)   {$\{d,i\}$};
  \node[scc, fill=blue!10]   (s3) at (1.5,-1.2){$\{c,h\}$};
  \node[scc, fill=yellow!30] (s4) at (3,-1.2)  {$\{e,j\}$};
  \draw[->, thick] (s1) -- (s2);
  \draw[->, thick] (s1) -- (s3);
  \draw[->, thick] (s2) -- (s4);
  \draw[->, thick] (s3) -- (s4);
  \node[below=4pt, keyred, font=\small] at (1.5, -2) {$G^{SCC}$ (DAG)};
\end{tikzpicture}
\end{center}

% ══════════════════════════════════════════════════════════════════════════════
\section*{Flow Networks}
\addcontentsline{toc}{subsection}{Flow Networks}
% ══════════════════════════════════════════════════════════════════════════════

% ─────────────────────────────────────────────────────────────────────────────
\subsection{Definition}

\begin{definition}{Flow network}{flownet}
A \textbf{flow network} is a directed graph $G = (V, E)$ where:
\begin{itemize}[noitemsep]
  \item Each edge $(u,v)$ has a \textbf{capacity} $c(u,v) \ge 0$
        (and $c(u,v) = 0$ if $(u,v) \notin E$).
  \item There is a \textbf{source} $s$ and a \textbf{sink} $t$.
  \item We assume the absence of anti-parallel edges (without loss of generality).
\end{itemize}
\end{definition}

\textbf{Example:} transporting Gruyère cheese ($s$) to Lausanne ($t$) via
intermediate cities, each route having a maximum capacity.

\begin{center}
\begin{tikzpicture}[font=\small, >=Stealth,
    node/.style={circle, draw, fill=blue!8, minimum size=0.7cm, inner sep=1pt}]
  \node[node, fill=green!20] (s)  at (0, 1)   {$s$};
  \node[node]                (v1) at (2, 2)   {$v_1$};
  \node[node]                (v2) at (2, 0)   {$v_2$};
  \node[node]                (v3) at (4, 2)   {$v_3$};
  \node[node]                (v4) at (4, 0)   {$v_4$};
  \node[node, fill=red!20]   (t)  at (6, 1)   {$t$};
  \draw[->, thick] (s)  -- node[above, font=\tiny]{3} (v1);
  \draw[->, thick] (s)  -- node[below, font=\tiny]{2} (v2);
  \draw[->, thick] (v1) -- node[above, font=\tiny]{2} (v3);
  \draw[->, thick] (v1) -- node[right, font=\tiny]{1} (v2);
  \draw[->, thick] (v2) -- node[below, font=\tiny]{3} (v4);
  \draw[->, thick] (v2) -- node[above, font=\tiny]{3} (v3);
  \draw[->, thick] (v3) -- node[above, font=\tiny]{3} (t);
  \draw[->, thick] (v4) -- node[below, font=\tiny]{2} (t);
  \draw[->, thick] (v3) -- node[right, font=\tiny]{1} (v4);
\end{tikzpicture}
\end{center}

\subsubsection{Removing anti-parallel edges}

If the graph contains $(u,v)$ and $(v,u)$, we replace $(u,v)$ with two edges
$(u, v')$ and $(v', v)$ with $c(u,v') = c(v',v) = c(u,v)$, by introducing a
dummy vertex $v'$. Repeat for each anti-parallel pair ($O(E)$ times).

% ─────────────────────────────────────────────────────────────────────────────
\subsection{Definition of a Flow}

\begin{definition}{Flow}{flow}
A \textbf{flow} is a function $f : V \times V \to \mathbb{R}$ satisfying:

\textbf{Capacity constraint:} $\forall\, u, v \in V$,
\[
  0 \le f(u,v) \le c(u,v).
\]

\textbf{Flow conservation:} $\forall\, u \in V \setminus \{s, t\}$,
\[
  \underbrace{\sum_{v \in V} f(v,u)}_{\text{incoming flow}} =
  \underbrace{\sum_{v \in V} f(u,v)}_{\text{outgoing flow}}.
\]
\end{definition}

The \textbf{value} of flow $f$ is:
\[
  |f| = \sum_{v \in V} f(s,v) - \sum_{v \in V} f(v,s)
       = \text{total flow exiting the source}.
\]

\textbf{Example} (network above, flow with value $|f| = 4$):

\begin{center}
\begin{tikzpicture}[font=\small, >=Stealth,
    node/.style={circle, draw, fill=blue!8, minimum size=0.7cm, inner sep=1pt}]
  \node[node, fill=green!20] (s)  at (0, 1)   {$s$};
  \node[node]                (v1) at (2, 2)   {$v_1$};
  \node[node]                (v2) at (2, 0)   {$v_2$};
  \node[node]                (v3) at (4, 2)   {$v_3$};
  \node[node]                (v4) at (4, 0)   {$v_4$};
  \node[node, fill=red!20]   (t)  at (6, 1)   {$t$};
  \draw[->, thick, keyblue] (s)  -- node[above, font=\tiny]{2/3} (v1);
  \draw[->, thick, keyblue] (s)  -- node[below, font=\tiny]{2/2} (v2);
  \draw[->, thick, keyblue] (v1) -- node[above, font=\tiny]{2/2} (v3);
  \draw[->, thick, keyblue] (v1) -- node[right, font=\tiny]{0/1} (v2);
  \draw[->, thick, keyblue] (v2) -- node[below, font=\tiny]{1/3} (v4);
  \draw[->, thick, keyblue] (v2) -- node[above, font=\tiny]{1/3} (v3);
  \draw[->, thick, keyblue] (v3) -- node[above, font=\tiny]{3/3} (t);
  \draw[->, thick, keyblue] (v4) -- node[below, font=\tiny]{1/2} (t);
  \draw[->, thick, keyblue] (v3) -- node[right, font=\tiny]{0/1} (v4);
  \node[below=4pt, font=\small] at (3, -0.8) {Flow $f/c$, value $|f|=4$};
\end{tikzpicture}
\end{center}

% ─────────────────────────────────────────────────────────────────────────────
\subsection{Residual Network}

\begin{definition}{Residual Network $G_f$}{residual}
Given a flow $f$ in $G = (V, E)$, the \textbf{residual capacity} is:
\[
  c_f(u,v) =
  \begin{cases}
    c(u,v) - f(u,v) & \text{if } (u,v) \in E \quad (\text{remaining capacity})\\
    f(v,u)          & \text{if } (v,u) \in E \quad (\text{cancellable flow}) \\
    0               & \text{otherwise.}
  \end{cases}
\]
The \textbf{residual network} is $G_f = (V, E_f)$ where
$E_f = \{(u,v) \in V \times V : c_f(u,v) > 0\}$.
\end{definition}

\begin{remark}{Role of reversed residual edges}{residualinv}
The reversed edges $(v,u)$ in $G_f$ represent the possibility of
\emph{reducing} the flow on $(u,v)$, which is equivalent to rerouting it elsewhere.
It is this mechanism that allows the algorithm to correct bad initial choices.
\end{remark}

% ─────────────────────────────────────────────────────────────────────────────
\subsection{Ford-Fulkerson Method}

\begin{algorithm}[H]
\caption{Ford-Fulkerson-Method($G$, $s$, $t$)}
\begin{algorithmic}[1]
\State Initialize $f(u,v) \leftarrow 0$ for all $(u,v)$
\While{there exists an \textbf{augmenting path} $p$ from $s$ to $t$ in $G_f$}
    \State \textbf{Augment} the flow $f$ along $p$
\EndWhile
\State \Return $f$
\end{algorithmic}
\end{algorithm}

An \textbf{augmenting path} is a simple path from $s$ to $t$ in the residual
network $G_f$. We send a flow equal to the \textbf{minimum residual capacity}
on this path:
\[
  c_f(p) = \min_{(u,v) \in p} c_f(u,v).
\]

For each edge $(u,v) \in p$, the flow is updated by:
$f(u,v) \leftarrow f(u,v) + c_f(p) - c_f(v,u)$ (taking into account reversed 
edges).

\subsubsection{Step-by-step example}

Diamond graph with a central edge of capacity 1 (the two paths of capacity 1
each):

\begin{center}
\begin{tikzpicture}[font=\small, >=Stealth,
    node/.style={circle, draw, fill=blue!8, minimum size=0.6cm, inner sep=1pt}]
  % Step 1: initial flow 0
  \begin{scope}[xshift=0cm]
    \node[node, fill=green!20] (s1) at (0,1)  {$s$};
    \node[node]                (m1) at (1.4,1){};
    \node[node, fill=red!20]   (t1) at (2.8,1){$t$};
    \draw[->,thick] (s1) -- node[above,font=\tiny]{0/1} (m1);
    \draw[->,thick] (m1) -- node[above,font=\tiny]{0/1} (t1);
    \draw[->,thick] (s1) to[bend right=30] node[below,font=\tiny]{0/1} (m1);
    \draw[->,thick] (m1) to[bend right=30] node[below,font=\tiny]{0/1} (t1);
    \node[below=6pt,font=\tiny] at (1.4,0.4) {initial flow};
  \end{scope}
  % Step 2: after 1st augmenting path
  \begin{scope}[xshift=4cm]
    \node[node, fill=green!20] (s2) at (0,1)  {$s$};
    \node[node]                (m2) at (1.4,1){};
    \node[node, fill=red!20]   (t2) at (2.8,1){$t$};
    \draw[->,thick,keygreen] (s2) -- node[above,font=\tiny]{1/1} (m2);
    \draw[->,thick,keygreen] (m2) -- node[above,font=\tiny]{1/1} (t2);
    \draw[->,thick] (s2) to[bend right=30] node[below,font=\tiny]{0/1} (m2);
    \draw[->,thick] (m2) to[bend right=30] node[below,font=\tiny]{0/1} (t2);
    \node[below=6pt,font=\tiny] at (1.4,0.4) {$|f|=1$};
  \end{scope}
  % Step 3: lower path, flow=2
  \begin{scope}[xshift=8cm]
    \node[node, fill=green!20] (s3) at (0,1)  {$s$};
    \node[node]                (m3) at (1.4,1){};
    \node[node, fill=red!20]   (t3) at (2.8,1){$t$};
    \draw[->,thick,keygreen] (s3) -- node[above,font=\tiny]{1/1} (m3);
    \draw[->,thick,keygreen] (m3) -- node[above,font=\tiny]{1/1} (t3);
    \draw[->,thick,keyblue] (s3) to[bend right=30] node[below,font=\tiny]{1/1} (m3);
    \draw[->,thick,keyblue] (m3) to[bend right=30] node[below,font=\tiny]{1/1} (t3);
    \node[below=6pt,font=\tiny] at (1.4,0.4) {$|f|=2$ (optimal)};
  \end{scope}
\end{tikzpicture}
\end{center}

\textbf{Case where reversed edges are necessary.} Without them, a bad
first choice (path passing through the central edge) would stall at $|f|=1$ instead
of 2. The residual network allows "undoing" this choice via a back edge.

% ─────────────────────────────────────────────────────────────────────────────
\subsection{Lecture Summary}

\begin{tcolorbox}[colback=lightblue, colframe=keyblue, fonttitle=\bfseries,
                  title={Key Points}]
\begin{itemize}[noitemsep]
  \item \textbf{SCC:} maximal set of mutually reachable vertices.
  \item \textbf{$G^{SCC}$ is a DAG.}
  \item \textbf{Kosaraju:} 2 DFS ($G$ then $G^T$ in decreasing order of $u.f$).
        Complexity $\Theta(V+E)$.
  \item \textbf{Flow Network:} capacity constraint $0 \le f(u,v) \le c(u,v)$
        and flow conservation at every vertex $\ne s, t$.
  \item \textbf{Residual Network $G_f$:} encodes what can still be sent
        (or cancelled) on each edge.
  \item \textbf{Ford-Fulkerson:} augment along $s \leadsto t$ paths in
        $G_f$ until exhaustion.
\end{itemize}
\end{tcolorbox}